\section{Theory: Four Mechanisms from International Relations}\label{sec:theory}

The bargaining literature in international relations is a natural home for the trade-request problem. It studies how parties with conflicting interests settle disputes when commitments cannot be enforced by a third party, when private information gives each side a reason to misrepresent its resolve, and when public statements change what each side can later accept. All three conditions hold here. I draw on four mechanisms.

\subsection{The commitment problem}

\citet{fearon1995} identifies commitment problems as one of three rationalist explanations for costly bargaining failure. Two parties may prefer a negotiated settlement to conflict, yet fail to reach one because one side cannot credibly promise not to exploit a future shift in power. \citet{powell2006} generalises the point: large, rapid shifts in the distribution of power make today's bargains worthless tomorrow.

A long-term NHL contract has the same structure. The team commits money, cap space and roster construction for six, seven or eight years. In return, the player promises his prime seasons. The team's side of the bargain is legally enforceable, because the salary is guaranteed and the clause binds. The player's side is not. Labour law and the CBA guarantee that he will be paid, but nothing guarantees that he will want to stay, and the things that matter most to the team (effort, leadership and willingness to sign the next contract) cannot be written into a contract. In the language of incomplete contracts \citep{hartmoore1988,hartmoore1990}, the player retains residual control over his own engagement. The player's commitment is therefore not \emph{self-enforcing}. When his outside option improves or his value of staying falls (a rebuild, a coaching change, a family move), the bargain is renegotiated in fact even though it is not renegotiated in law.

This is what the paper means by \emph{defection from the implicit commitment}. A trade request does not breach the NTC or NMC. The clause is a right the player holds, and waiving it is entirely his choice. What breaks is the unwritten half of the long-term deal.

\subsection{Tying hands, turned around}

\citet{schelling1960} observed that the power to bind oneself is the power to constrain an adversary. A negotiator who can credibly eliminate his own options forces the other side to accommodate. \citet{fearon1997} formalises this as \emph{tying hands}: a signal is credible if sending it and then backing down is costly.

In the standard story, the clause ties the team's hands for the player's benefit: the team cannot move him. The twist here is that a trade request reverses the direction. Once the player announces that he wants out, the team's hands are tied in the market for his services. Every rival general manager now knows that the incumbent must sell, so the incumbent's implicit threat to keep the player loses credibility. The clause that was meant to protect the player against the team becomes, in the hands of a player who wants to leave, a weapon against the team.

\subsection{Veto players and agenda control}

\citet{tsebelis2002} shows that adding a veto player shrinks the set of feasible policy changes (the ``winset''). A player with a full NMC is a veto player over every possible transaction. When he demands a trade, he does not merely accept or reject offers. He sets the agenda by naming the teams he will accept. That narrows the field of buyers, sometimes to one. The NHL's own language reflects this: Ottawa's general manager said Brady Tkachuk's clause ``changes how the transaction works'' and that the process was ``pointed at one team at the end'' \citep{cp24_tkachuk}. In auction terms, the veto converts a competitive sale into a bilateral negotiation, and \citet{bulowklemperer1996} show that competition among buyers is worth more to a seller than almost any amount of negotiating skill. The surplus the seller loses does not go to the player. It goes to the favoured buyer.

\subsection{Audience costs}

\citet{fearon1994} argues that leaders who make public threats incur \emph{audience costs} if they back down, and that these costs make public threats more credible than private ones. \citet{tomz2007} provides experimental evidence for such costs, while \citet{snyderborghard2011} question how often they bind. A public trade request is a public threat. It raises the player's own cost of reconciling with the team (fans, teammates and media will treat a reversal as weakness or insincerity) and therefore makes his exit credible. It also destroys the team's ability to bluff. Privately, a team can tell suitors that it is happy to keep the player and see whether the situation improves. Once the request is public, that claim is no longer believable. Connor Hellebuyck's request, private for months and then confirmed publicly in August 2026, is the clearest recent example \citep{thw_larkin_hellebuyck,br_hellebuyck}.

\subsection{Holdup and sunk costs}

Two further ideas complete the toolkit. First, the trade-request problem is a form of the \emph{holdup} studied in transaction-cost economics \citep{kca1978,williamson1979,williamson1985}. The team makes relationship-specific investments (cap commitments, roster construction, a marketing identity built around the player) whose value depends on the player's continued cooperation, and the clause lets that quasi-rent be appropriated ex post. The holdup here is unusual because the party that appropriates the rent is not the counterparty: the player gets his preferred destination, while the money goes to the acquiring team. Second, decision makers are prone to honour sunk costs \citep{arkesblumer1985,staw1976}, and sports organisations are no exception \citep{stawhoang1995,camererweber1999}. A team that has just committed \$70 million to a captain may wait in the hope that he changes his mind, even when the expected value of waiting is negative. Detroit's handling of Dylan Larkin's 2026 request illustrates the temptation.
